% ----------------------------------------------------------
\chapter{Detalhamento do Alinhamento Selecionado}
% ----------------------------------------------------------

\section{Cálculo da Correção dos Azimutes}

Para a determinação da correção dos azimutes, considerou-se a declinação magnética inicial, a variação anual e a convergência meridiana, conforme dados disponíveis nas ortografias do Geoportal.



\textbf{Dados iniciais:}
\begin{itemize}
	\item Declinação magnética inicial: $D_i = -21^\circ 01'24''$ (em 01/03/2010);
	\item Variação anual: $v = -5{,}07'$ por ano;
	\item Convergência meridiana: $\gamma = 0^\circ 43'32{,}36''$;
	\item Data do projeto: 01/10/2027.
\end{itemize}

\subsection*{Período entre as épocas}
O intervalo de tempo entre as duas datas é:
\[
T = 17{,}5852\ \text{anos}.
\]

\subsection*{Variação total da declinação}
A variação acumulada da declinação magnética é:
\[
V_{\text{total}} = v \times T = -5{,}07' \times 17{,}5852 \approx -89{,}157' .
\]
Convertendo para graus, minutos e segundos:
\[
-89{,}157' = -1^\circ 29'9{,}42''.
\]

\subsection*{Declinação magnética final}
Somando a variação total à declinação inicial, obtém-se:
\[
D_f = D_i + V_{\text{total}} = -21^\circ 01'24'' - 1^\circ 29'9{,}42'' = -22^\circ 30'33{,}42'' \approx -22^\circ 30'33''.
\]

\subsection*{Relação entre Norte Magnético e Norte da Quadrícula}
Para determinar a diferença entre o Norte Magnético e o Norte da Quadrícula, adiciona-se a convergência meridiana à declinação magnética final:
\[
\Delta = D_f + \gamma = -22^\circ 30'33{,}42'' + 0^\circ 43'32{,}36'' \approx -21^\circ 47'01''.
\]

\subsection*{Correção a aplicar}
Portanto, a correção angular a ser aplicada entre o Norte Magnético e o Norte da Quadrícula, referente à data do projeto (01/10/2027), é:
\[
\boxed{-21^\circ 47'01''}.
\]

Esse valor deve ser somado (algebricamente) aos azimutes obtidos com base no Norte da Quadrícula, de modo a convertê-los em azimutes referidos ao Norte Magnético. Ressalta-se que todos os azimutes apresentados foram calculados em relação ao Norte da Quadrícula.









\begin{figure}[h]
	
	\includegraphics[width=0.85\textwidth]{unbtex-example/figuras/13.declinacao.png}
	\caption{Relação entre os nortes e declinações magnéticas.}
	\label{fig:declinacao}
\end{figure}











% ----------------------------------------------------------
\section{Dados Gerais do Alinhamento}
% ----------------------------------------------------------

Os dados gerais referentes ao alinhamento em estudo encontram-se na Tabela~\ref{tab:tab_3}.

\begin{table}[h]
	
	\caption{Dados gerais do alinhamento B.}
	\label{tab:tab_3}
	
		\begin{tabular}{|l|cc|cc|cc|c|c|}
			\hline
			\textbf{Ponto} & \multicolumn{2}{c|}{\textbf{Coord. Absolutas}} & 
			\multicolumn{2}{c|}{\textbf{Coord. Relativas}} & 
			\multicolumn{2}{c|}{\textbf{Comprimento}} & 
			\textbf{Deflexão} & \textbf{Azimute} \\ \cline{1-7}
			& X & Y & Xr & Yr & Calculado & Medido & (G,M,S) & (G,M,S) \\ \hline
			\textbf{PP} & 213966.9 & 8243813 & - & - & - & - & - & 159°02'24'' \\ \hline
			\textbf{PI-1} & 214289.1 & 8242972 & 322.2153 & -841.16 & 900.762 & 900.762 & 31°47'43'' & 127°14'41'' \\ \hline
			\textbf{PI-2} & 215304.1 & 8242200 & 1015.001 & -771.6787 & 1275.035 & 1275.035 & 24°43'53'' & 151°58'34'' \\ \hline
			\textbf{PI-3} & 215798.8 & 8241271 & 494.7447 & -929.5503 & 1053.013 & 1053.013 & 16°46'49'' & 168°45'24'' \\ \hline
			\textbf{PF} & 216147.9 & 8239515 & 349.0293 & -1755.771 & 1790.127 & 1790.127 & - & - \\ \hline
	\end{tabular}
\end{table}

O detalhamento gráfico do alinhamento está apresentado na Figura~\ref{fig:detalhamento_B}.

\begin{figure}[h]
	
	\includegraphics[width=0.85\textwidth]{unbtex-example/figuras/12.detalhamento_tracado.pdf}
	\caption{Detalhamento do traçado B.}
	\label{fig:detalhamento_B}
\end{figure}




